\section{Identification and Estimation}\label{sec:identification}

Three objects require estimation: the type-specific cost distributions $F_c^k$, the auction-level scale $a_t$, and the equilibrium entry counts. The third is directly observed. The first two require structure---the first by design of the Preg\~ao auction, the second through the variance decomposition in Assumption~\ref{a:uh}.

\subsection{Type-specific cost distributions}

Under Assumption~\ref{a:ipv} and the English-reverse interpretation, the drop-out price of each loser is a point observation of that loser's cost. The empirical distribution of the normalized drop-outs within a stratum is therefore an unbiased estimator of $F_c^k$---losers-only, it would seem, is all one needs. Losers-only, however, omits the lowest cost draw per auction (the winner's), biasing $F_c^k$ systematically upward at the left tail. I therefore use an \emph{all-bidders} ECDF, which adds the winner's final clock price to the stratum-specific ECDF. In a descending-clock auction this final price is approximately $c_{(2)}$, the second-to-last exit---an upper bound on the winner's cost but a substantially tighter one than the zero imputed by losers-only omission. The bias in all-bidders runs in the opposite direction and is materially smaller: losers-only overstates $F_c^k$ in the left tail; all-bidders understates it slightly throughout.

A Turnbull NPMLE, reported in Section~\ref{sec:robustness}, treats the winner formally as left-censored at $c_{(2)}$ and delivers point-identified $F_c^k$ without either bias. Across the eight (class $\times$ type $\times$ period) strata (Table~\ref{tab:v3_turnbull_fc}), the median estimates satisfy Turnbull $\le$ all-bidders $\le$ losers-only in all 8 strata, and the 75th-percentile estimates satisfy all-bidders $\le$ Turnbull $\le$ losers-only in 6 of 8 strata (the two pharmaceutical non-SME cells reverse, with Turnbull marginally below all-bidders by 0.01--0.02 of the reference price). The BNE decomposition computed on the Turnbull input differs from the all-bidders baseline by at most 16 percent of the total effect across classes (Table~\ref{tab:v3_sensitivity_fc}), with the largest gap in pharmaceutical Turnbull. All-bidders is the main-text estimator; Turnbull is the robustness benchmark; losers-only is the input the DiD benchmark of Appendix~\ref{app:did} effectively absorbs through the item fixed effects.

A cross-modality consistency check supports the descending-clock specification. Convite auctions in the same sample are first-price sealed-bid; the \citet{gpv2000} inversion recovers $F_c^k$ from the Convite bid distribution. The Convite GPV and Preg\~ao drop-out estimates coincide to approximately one percentage point at the median in pharma non-SME Pre---the stratum where the comparison is least confounded by the policy reform, since the period predates the cutoff and the bidders are non-SMEs unaffected by the set-aside. In other strata the gap is wider and motivates the auction-level heterogeneity correction below. Figure~\ref{fig:v3_cross_modality_uh} plots the two recoveries side by side for all four strata. The check also serves as an indirect collusion screen: the bidder-group detection tests of \citet{conley2016} would flag systematic divergence between GPV-inverted bids and drop-out exits in collusive settings, and the convergence reported here is consistent with non-collusive equilibrium play in the stratum tested. A formal test of $F_c^{\text{SME}} = F_c^{\neg}$ (symmetric IPV against the asymmetric reading) is not the central exercise, but the recovered cell-level medians and 75th-percentile estimates differ by 0.10--0.30 of the reference price across the (pharma $\times$ type) cells (Table~\ref{tab:v3_pregao_fc_summary}), well outside Monte Carlo sampling noise; asymmetry is therefore consistent with the data, not just imposed by Assumption~\ref{a:ipv}.

\begin{figure}[!htbp]
\centering
\includegraphics[width=0.98\textwidth]{../output/figures/fig_v3_cross_modality_uh.pdf}
\caption[Cross-modality identification check after UH correction.]{\textit{Cross-modality identification check, UH-clean $F_{c_\epsilon}$.} Each panel plots the type-specific cost distribution estimated two ways on the same sample: the \citet{gpv2000} inversion on Convite first-price sealed-bid auctions (solid) and the \citet{haile2003} drop-out point-ID on Preg\~ao descending-clock auctions (dashed). Pre and Post periods are colour-coded. Convergence in the pharma non-SME Pre panel (upper right) supports the IPV restriction in that stratum: two independent structural recoveries on the same primitive land within sampling error of each other. Residual gaps in other panels motivate the auction-level heterogeneity correction in Section~\ref{sec:identification}.}
\label{fig:v3_cross_modality_uh}
\end{figure}

\subsection{Auction-level heterogeneity}

The median-cost gap across modalities in non-pharma---up to 23 percentage points of the reference price---is too large to attribute to sampling noise. Reference-price normalization absorbs most of the scale variation across items, but not all of it: two items in the same CADMAT class can differ in ANVISA registration burden, in contract complexity, or in volume-specific cost elasticities in ways the reference price does not fully reflect. The residual common component is Assumption~\ref{a:uh}'s $a_t$.

I estimate the variance decomposition by method of moments. Within each stratum $(k, \text{pharma}, \text{period})$, $\sigma_a^{2,k}$ is the between-auction variance of auction-mean log-bids, and $\sigma_e^{2,k}$ is the residual within-auction variance. The intraclass correlation
\[
\rho^k = \frac{\sigma_a^{2,k}}{\sigma_a^{2,k} + \sigma_e^{2,k}}
\]
ranges across the four Preg\~ao (pharma $\times$ type) cells from approximately 0.28 (Preg\~ao pharmaceutical SMEs, average across periods, with a low-$n$ Pre-period stratum pulling the cell mean down) through 0.36 (Preg\~ao non-pharmaceutical SMEs) and 0.39 (Preg\~ao non-pharmaceutical non-SMEs) up to approximately 0.59 (Preg\~ao pharmaceutical non-SMEs); the corresponding Convite cells deliver 0.55 and 0.66 (Table~\ref{tab:v3_uh_variance}). Restricted to the three Preg\~ao cells outside the low-$n$ pharma-SME-Pre stratum (range 0.36--0.59), the order of magnitude is comparable to the 0.50--0.70 range Krasnokutskaya (2011) reports for California highway contracts. A best-linear-predictor shrinkage then estimates each auction's specific heterogeneity as
\[
\hat{a}_t = \frac{\hat\sigma_a^2 \, N_t}{\hat\sigma_a^2 \, N_t + \hat\sigma_e^2} \cdot \left( \bar y_t - \bar y \right),
\]
and the UH-clean residuals $\hat e_{it} = y_{it} - \hat a_t$ deliver cleaned cost estimates $c_\epsilon^{\text{clean}} = \exp(\hat e_{it})$. These cleaned bids are the inputs to all downstream structural estimation.

The correction closes ninety-three percent of the cross-modality median gap in pharma SME Post, twenty-three percent in pharma non-SME Post, and leaves the non-pharma gap largely intact---evidence that auction-level heterogeneity is the binding constraint in pharma but only part of the story in non-pharma, where richer item-level heterogeneity presumably operates (and is addressed partially by the sample filter $c_\epsilon \le 3$). The Krasnokutskaya shrinkage assumes Gaussianity of $a_t$ and $e_{it}$; the Turnbull NPMLE in Section~\ref{sec:robustness} dispenses with this assumption and ratifies the estimates.

\subsection{Entry}

Entry is modeled as Poisson arrivals with stratum-specific rates. I take the observed Pre-period average counts of SME and non-SME bidders per auction as the status-quo arrival rates, and the Post-period SME counts as the equilibrium arrival under the restricted regime. No structural entry cost is jointly estimated with the cost primitives; this is the reduced-form treatment in the \citet{atheyseira2011} spirit. The same logic applies to the post-policy SME bidder distribution: in the main specification, $F_c^{\text{SME,Post}}$ is treated as an equilibrium-selection object generated by the protected regime rather than forced to equal $F_c^{\text{SME,Pre}}$. Section~\ref{sec:robustness} reports the stricter limiting case in which that equality is imposed as a benchmark. A zero-profit calculation applied to the simulated winning-margin profit distribution delivers an implied entry cost per bidder in the R\$0.11--R\$2.46 range (Table~\ref{tab:v3_entry_cost}), with non-SMEs facing roughly 4.5 times the SME entry cost across both classes ($\kappa^{\neg}/\kappa^{\text{SME}} = 4.5$ in non-pharmaceuticals and $4.6$ in pharmaceuticals). These magnitudes are plausible given the documentation and certification burden of responding to a BEC tender.

\subsection{Counterfactual simulation}

With $F_c^k$ and entry counts in hand, the BNE counterfactual is a Monte Carlo over draws from the estimated primitives. For each of $B = 2{,}000$ simulated auctions, I draw $n^{\text{SME}} \sim \mathrm{Poisson}(\bar N^{\text{SME}})$ and analogously for non-SMEs (the Poisson rate is the observed mean entry count by stratum, distinct from the MCPF $\lambda$ used in the welfare arithmetic of Section~\ref{sec:welfare}), sample costs from the corresponding $F_c^k$, sort, and record $c_{(1)}$ (the winner's cost, the production cost), $c_{(2)}$ (the expected transaction price under revenue equivalence), and the winner's type. Aggregating over draws yields the means, quantiles, and decompositions reported in Section~\ref{sec:results}. Standard errors come from a cluster bootstrap at the auction level with $B_{bs} = 500$ replicates; details are in Section~\ref{sec:robustness}.
